\section*{第 \ref{ch:g12:domesticated-plants} 章——驯化植物}

\begin{solution}{exo:g12:domesticated-plants:1}
\omterm{def:g12:domesticated-plants:domestication}{驯化}：野生\omterm{def:g10:biodiversity-scales:species}{物种}转化为依赖人类并为人类服务的\omterm{def:g10:biodiversity-scales:species}{物种}。
\omterm{def:g12:domesticated-plants:domestication}{人工选择}：人类保留并播种
具有所需性状的个体，以提高
与之相关的\omterm{def:g10:universal-dna:gene}{等位基因}的频率。
\end{solution}

\begin{solution}{exo:g12:domesticated-plants:2}
穗不易落粒（无种子扩散）；种子立即萌发（不必
等待糟糕的年份过去）；失去毒素（能被所有动物取食）；
无分枝植株上只有一个大穗（把所有鸡蛋放在一个篮子里）；
同时成熟（一场糟糕的天气就能毁掉整个作物）。
\end{solution}

\begin{solution}{exo:g12:domesticated-plants:3}
由两个纯合近交系杂交得到的后代，整齐、生长旺盛。
其自身的种子会分离出长势较弱的植株混杂群体，
因此农民每年都需购买新的杂交种子。
\end{solution}

\begin{solution}{exo:g12:domesticated-plants:4}
将\omterm{def:g10:universal-dna:gene}{基因}插入土壤细菌天然转入植物\omterm{def:g10:cells-common-unit:cell}{细胞}的 \omterm{def:g10:universal-dna:information}{DNA} 中；
获得该\omterm{def:g10:universal-dna:gene}{基因}的受感染\omterm{def:g10:cells-common-unit:cell}{细胞}
再生为完整植株。
\end{solution}

\begin{solution}{exo:g12:domesticated-plants:5}
由约 \num{1.9} 增至 \qty{7.2}{t/ha}，近四倍。
\end{solution}

\begin{solution}{exo:g12:domesticated-plants:6}
籽粒易落粒的植株在收获前或收获过程中就损失了种子；
茎秆坚韧的\omterm{def:g11:mutations:mutation}{突变}体则保住了种子，被采收下来，
因而也被留作播种。收获本身逐代选择了
坚韧茎秆的\omterm{def:g10:universal-dna:gene}{等位基因}。
\end{solution}

\begin{solution}{exo:g12:domesticated-plants:7}
同一\omterm{def:g10:biodiversity-scales:species}{物种}：它们能够杂交并产生可育后代。农民在各
谱系中分别选择了不同器官的过度发育——
叶片、顶芽、花芽——因此它们的差异在于生长的调控，
而非杂交能力。
\end{solution}

\begin{solution}{exo:g12:domesticated-plants:8}
F1 全部为 $Aa$。F2 中：$AA$ $1/4$，$Aa$ $1/2$，$aa$ $1/4$。F1 整齐一致，
是因为每株植株从一个亲本获得 $A$，从另一亲本获得 $a$；
F2 中\omterm{def:g12:meiosis-genetic-shuffling:meiosis}{减数分裂}与受精使\omterm{def:g10:universal-dna:gene}{等位基因}
随机分配。
\end{solution}

\begin{solution}{exo:g12:domesticated-plants:9}
数以百万计的螟虫中，只有少数携带使其耐受毒素的\omterm{def:g11:mutations:mutation}{突变}；
在每株植物都产生毒素的田地上，只有它们存活
并繁殖——此即选择。避难所使大量敏感螟虫得以存活；
它们与少数抗性个体交配，产生携带一个抗性\omterm{def:g10:universal-dna:gene}{等位基因}的后代；
此类后代因该\omterm{def:g10:universal-dna:gene}{等位基因}为\omterm{def:g11:genes-and-disease:genetic}{隐性}而在含毒玉米上死亡，
因此该\omterm{def:g10:universal-dna:gene}{等位基因}每代都被稀释。
\end{solution}

\begin{solution}{exo:g12:domesticated-plants:10}
相关\omterm{def:g10:universal-dna:gene}{基因}控制发育——如植株分枝、
籽粒外壳的生长、行数等；某一\omterm{def:g10:universal-dna:gene}{基因}调控上的改变
会改变整个器官的形态。
调控上的少数改变，带来显著的可见效应。
\end{solution}

\begin{solution}{exo:g12:domesticated-plants:11}
它们保存着选择从作物中剔除的\omterm{def:g10:universal-dna:gene}{等位基因}：对病害和害虫的抗性，
以及对干旱、盐碱或低温的耐受性。当新的病害或气候出现时，
育种者将这些\omterm{def:g10:universal-dna:gene}{等位基因}重新导入作物；
一旦野生近缘种灭绝，其\omterm{def:g10:universal-dna:gene}{等位基因}便随之消失。
\end{solution}

\begin{solution}{exo:g12:domesticated-plants:12}
油菜的花粉使芥菜受精，杂交种携带耐受\omterm{def:g10:universal-dna:gene}{基因}；
经若干代回交芥菜，并在田间喷施除草剂的选择下，
耐受型野生芥菜得以扩散。过程包括杂交，
以及\omterm{def:g10:universal-dna:gene}{基因}在\omterm{def:g12:selection-drift-speciation:population}{种群}间的水平转移，
随后是除草剂导致的\omterm{def:g12:selection-drift-speciation:selection}{自然选择}。
\end{solution}

\begin{solution}{exo:g12:domesticated-plants:13}
生物学方面：植物的\omterm{def:g10:cell-metabolism:autotrophy}{光合作用}及籽粒充实能力
已接近极限，而进一步增产所依赖的优良品种
\omterm{def:g10:biodiversity-scales:biodiversity}{遗传多样性}狭窄。
非生物学方面：化肥和农药的使用量不再增加，
正在减少，且气候胁迫日益加重。育种者可以
从野生近缘种引入\omterm{def:g10:universal-dna:gene}{等位基因}，或编辑\omterm{def:g10:cell-metabolism:autotrophy}{光合作用}
及耐旱性相关的\omterm{def:g10:universal-dna:gene}{基因}。
\end{solution}

\begin{solution}{exo:g12:domesticated-plants:14}
杂交育种：所需\omterm{def:g10:universal-dna:gene}{等位基因}须存在于可与作物杂交的植株中；
数以千计的\omterm{def:g10:universal-dna:gene}{基因}被重新组合，须经约十代选择加以筛选；
\omterm{prop:g10:universal-dna:universal}{转基因技术}：\omterm{def:g10:universal-dna:gene}{基因}可来自任何生物；只增加一个\omterm{def:g10:universal-dna:gene}{基因}，
其余\omterm{def:g10:universal-dna:gene}{基因}组保持不变；
仅需数年。
\end{solution}

\begin{solution}{exo:g12:domesticated-plants:15}
两者的算法完全相同：携带某些\omterm{def:g10:universal-dna:gene}{等位基因}的个体繁殖得更多，
\omterm{def:g10:universal-dna:gene}{等位基因}的频率随之上升。二者的不同在于作用因素——以人类的选择取代环境，
以及对植物而言的结果：\omterm{def:g12:selection-drift-speciation:selection}{自然选择}保留能够独立生存的\omterm{def:g10:biodiversity-scales:species}{物种}，
\omterm{def:g12:domesticated-plants:domestication}{驯化}则选择那些服务于选择者的性状，
使植物依赖于人类。\omterm{def:g12:domesticated-plants:domestication}{驯化}是
由不同选择者施行的\omterm{def:g12:selection-drift-speciation:selection}{自然选择}，
而非其对立面。
\end{solution}

\begin{solution}{pb:g12:domesticated-plants:1}
\textbf{1.} 大刍草 $12 \times 0.08 \approx \qty{1}{g}$；玉米 $600
\times 0.3 = \qty{180}{g}$：接近两百倍。

\textbf{2.} 坚硬穗轴上籽粒众多（不利扩散）；裸粒
（无保护，易受动物和腐烂侵害）；单茎上只长一个大果穗
（风险不分散）；籽粒包于苞叶之中（
完全无法落到地面上）。

\textbf{3.} 裸粒\omterm{def:g10:universal-dna:gene}{等位基因}为\omterm{def:g11:genes-and-disease:genetic}{隐性}：植株只有携带两个拷贝时才表现为裸粒，
而对于稀有\omterm{def:g10:universal-dna:gene}{等位基因}，这种情况发生在
极少数植株中；携带一个拷贝的植株外观正常。

\textbf{4.} $q^2 = 0.0001$：万株中有一株。

\textbf{5.} 所有种子均来自 $nn$ 植株，其籽粒从母本获得
$n$；花粉主要来自普通田块（$N$
频率为 0.99）。种子几乎全为 $Nn$，故
$n$ 的频率跃升至约 0.5——仅一代强选择。

\textbf{6.} 由 $600/12 = 50$，历经 9000 代，增长倍数为 $50^{1/9000}
\approx 1.0004$ 每代，增益为 0.04\%。
微小的步进经数千代累积，便成为巨大变化——
选择的算术。

\textbf{7.} 若众多\omterm{def:g10:universal-dna:gene}{基因}各自贡献微小，选择会提高所有这些\omterm{def:g10:universal-dna:gene}{基因}中有利\omterm{def:g10:universal-dna:gene}{等位基因}的
频率，并能长时间不断找到新的组合；
单\omterm{def:g10:universal-dna:gene}{基因}性状则在\omterm{def:g11:genes-and-disease:genetic}{隐性}\omterm{def:g10:universal-dna:gene}{等位基因}变为纯合时即
发生转换。

\textbf{8.} 同属一个\omterm{def:g10:biodiversity-scales:species}{物种}：杂种可育。\omterm{def:g12:domesticated-plants:domestication}{驯化}改变了
少数\omterm{def:g10:universal-dna:gene}{基因}上\omterm{def:g10:universal-dna:gene}{等位基因}的频率及植株形态，但并未
产生生殖隔离。

\textbf{9.} 能攻克某一株植物的病害或害虫，便能攻克全部
植株；答案在于老品种和野生近缘种中保存的多样性，
存于种子库与田间。

\textbf{10.} 农民既是植物的扩散者，也是其选择环境，
并且是其通向下一代的唯一途径：玉米
每一代都处于\omterm{def:g12:domesticated-plants:domestication}{人工选择}之下，无法脱离。

\textbf{11.} $10/5.5 - 1 \approx 82\%$，高于较优亲本；
$1 - 8.2/10 = 18\%$，低于 F1。

\textbf{12.} 每个品系在若干弱\omterm{def:g11:genes-and-disease:genetic}{隐性}\omterm{def:g10:universal-dna:gene}{等位基因}上纯合；
在 F1 中，每个此类\omterm{def:g10:universal-dna:gene}{等位基因}都被另一品系功能正常的\omterm{def:g10:universal-dna:gene}{等位基因}掩盖。
在 F2 中，四分之一的植株在每个此类\omterm{def:g10:universal-dna:gene}{基因}上
均为弱\omterm{def:g10:universal-dna:gene}{等位基因}纯合，平均值随之下降。

\textbf{13.} 购买：$10 - 0.4 = \qty{9.6}{t/ha}$；自留：\qty{8.2}{t/ha}。
购买每年多收 \qty{1.4}{t/ha}，这正是
掌握自交系的公司能够每季出售杂交种子的原因。

\textbf{14.} 公司保有自交系；若农民掌握它们，便可
无限期地自行制种，而公司只能一次性售出。
自交系是公司的资本；出售的只是其产物。

\textbf{15.} 杂交品种的每块田在遗传上完全一致，
各地只种植少数几个杂交品种，多样性随之缩窄。农民
不再自行制种，每年都依赖供种者。

\textbf{16.} $q^2 = 10^{-6}$：每百万只玉米螟中仅一只存活；在
存活者中，全部为 $rr$，$r$ 的频率为 1。

\textbf{17.} 存活者彼此交配，其后代
全部为 $rr$；数代之内，田间玉米螟
均已耐药——正如疗程不完整的患者体内的细菌一样。

\textbf{18.} 耐药存活者（$rr$）与避难区玉米螟
（$RR$，占绝大多数）交配，产生 $Rr$ 后代，它们在毒素玉米上死亡。
只要敏感玉米螟远多于耐药个体，
几乎每只耐药玉米螟的后代都是杂合的，并被
清除，因此 $r$ 始终稀少。

\textbf{19.} 相同之处：二者都旨在防止稀有
耐药个体的增殖。不同之处：完整疗程通过
在耐药性被选择之前杀死全部个体；避难区
则容纳存活者，并以敏感个体将其稀释。

\textbf{20.} 每穗籽粒多出近两百倍；
约五个主效\omterm{def:g10:universal-dna:gene}{基因}；一种依赖农民扩散和播种
自身的植物。
\end{solution}
